\section{Bellarmine and Galileo: The Documents and Their Limits
(1611--1616)}
\label{sec:galileo}

No part of Bellarmine's life is more written about, or more
overwritten, than his part in the Copernican affair of 1615--1616.
This section therefore proceeds document by document. Every text
quoted is cited from Antonio Favaro's national edition of Galileo's
works (\work{Le Opere di Galileo Galilei, Edizione Nazionale}),
with the volume, page, and the manuscript witness Favaro identifies.
English renderings, unless otherwise credited, are project
translations of Favaro's text and are marked as such. What the
documents do not say is stated as plainly as what they do.

\subsection{1611: the cardinal asks the astronomers}

When Galileo came to Rome in the spring of 1611 with his telescope,
Bellarmine did two documented things. He looked---``I too have
seen, through the same instrument, some very marvelous things
concerning the moon and Venus,'' he wrote---and he asked his own
experts. His letter of 19~April 1611 to the mathematicians of the
Roman College (Clavius among them) lists five findings with
precision---the crowds of invisible fixed stars and the composition
of the Milky Way; the strange shape of Saturn; the phases of Venus;
the rough and unequal surface of the moon; the four fast-moving
satellites of Jupiter---and asks the fathers ``to tell me sincerely
your opinion,'' because he hears various views (\work{Opere}~XI,
no.~515, from the original with autograph signature in the
Biblioteca Nazionale, Florence; A/N; project translation). Their
reply of 24~April confirmed the observations point by point, with
Clavius's reservation about whether the moon's surface itself was
uneven (no.~520; N). The exchange shows the man exactly: a
non-mathematician of open curiosity, consulting competence,
committed to nothing. A courteous note to Galileo of 23~June 1612,
thanking him for his book on floating bodies, closes the file of
their early relations (\work{Opere}~XI, no.~709; A).

\subsection{12 April 1615: the letter to Foscarini}

By 1615 the Carmelite provincial Paolo Antonio Foscarini had
published a letter arguing that the Copernican system was
compatible with Scripture, and Galileo's own Copernican letters
were circulating. Bellarmine's reply to Foscarini of 12~April
1615---which he plainly intended Galileo to read, naming him in
its first article---survives in a contemporary copy (Biblioteca
della R. Accademia dei Lincei, Cod.\ Volpicelliano~A, fols.\
159r--160r; \work{Opere}~XII, no.~1110, pp.~171--172; A at copy
level). Its three numbered points deserve exact statement, for
they are the closest thing to Bellarmine's considered position on
science and Scripture.

First: ``It seems to me that Your Reverence and Signor Galileo act
prudently in contenting yourselves with speaking \latin{ex
suppositione} and not absolutely, as I have always believed
Copernicus spoke''; to say that supposing the earth's motion the
appearances are saved better than by eccentrics and epicycles ``is
very well said and carries no danger; and this suffices for the
mathematician''; but to affirm the sun's centrality in reality
``is a very dangerous thing, likely not only to irritate all the
philosophers and scholastic theologians, but also to harm the
Holy Faith by rendering the Holy Scriptures false.''

Second: the Council of Trent prohibits interpreting Scripture
against the common consent of the holy Fathers, and all
commentators, Greek and Latin, expound the sun's motion literally;
nor will it do to say the matter is not one of faith, ``for if it
is not a matter of faith \latin{ex parte obiecti}, it is a matter
of faith \latin{ex parte dicentis}''---on the side of the speaker,
the Holy Spirit.

Third---the sentence on which three centuries of interpretation
hang: ``I say that when there shall be a true demonstration
(\latin{vera demostratione}) that the sun stands in the center of
the world and the earth in the third heaven\ldots\ then it will be
necessary to proceed with great consideration in explaining the
Scriptures that appear contrary, and rather to say that we do not
understand them than to say that what is demonstrated is false.
But I shall not believe that there is such a demonstration until
it is shown to me''; and demonstrating that the supposition saves
the appearances is not the same as demonstrating the reality; of
the second ``I have very great doubt, and in case of doubt one
must not abandon Holy Scripture as expounded by the holy
Fathers.'' He closes with Solomon and the shore that seems to
recede from the ship: in the case of sun and earth, no wise man
needs to correct the appearance, ``for he clearly experiences that
the earth stands still.''

\witnesslimit{Document and interpretation. The letter's text
establishes: a methodological tolerance of Copernicanism as
supposition; a theological bar against asserting it as physical
truth absent demonstration; an explicit, conditional concession
that demonstration would force reinterpretation of Scripture; and
a personal judgment that no demonstration existed or was likely.
Whether the third point makes Bellarmine a far-sighted
methodologist, or the second makes him the architect of the
coming error, or both, is interpretation; both readings have
distinguished sponsors, and the 1907 Catholic Encyclopedia
already read the letter as stating the rule that a solidly
demonstrated theory ``must'' govern scriptural interpretation.
This study asserts only the text.}

\subsection{February--March 1616: censure, direction, minute,
decree}

Events then passed out of the epistolary register. Galileo came
to Rome to press the Copernican cause in person. On 19~February
1616 the Holy Office sent two propositions extracted from
Galileo's positions to its consultor theologians: that the sun is
the immobile center of the world; that the earth is not the
center, but moves both wholly and with daily motion. On
24~February eleven qualifiers signed their censure: the first
proposition ``foolish and absurd in philosophy, and formally
heretical, inasmuch as it expressly contradicts the statements of
Holy Scripture in many places according to the propriety of the
words and according to the common exposition and sense of the
holy Fathers and doctors of theology''; the second earning ``the
same censure in philosophy, and, in regard to theological truth,
at least erroneous in faith'' (\work{Opere}~XIX, pp.~320--321,
from the signed original, Vat.\ car.\ 377r; N; project
translation). Bellarmine's name is not among the eleven; the
censure was an internal act of consultors, never published as
such.

The next day, Thursday 25~February, the pope acted on the
censure. The Holy Office minute records: the pope ``ordered the
Most Illustrious Lord Cardinal Bellarmine to summon Galileo
before him and warn him (\latin{moneat}) to abandon the said
opinion; and, should he refuse to obey, the Father Commissary,
before notary and witnesses, is to enjoin on him a precept
(\latin{praeceptum}) to abstain altogether from teaching or
defending doctrine and opinion of this kind, or treating of it;
and, should he not acquiesce, he is to be imprisoned''
(\work{Opere}~XIX, p.~321, car.\ 378v; N; project translation).
The structure is graduated: warning first; precept only on
refusal; prison only on defiance.

What happened the following morning, Friday 26~February 1616, in
Bellarmine's residence, is recorded in the most disputed document
of the affair---an unsigned minute in the trial dossier (car.\
378v--379r; \work{Opere}~XIX, pp.~321--322; N). It states that
Bellarmine warned Galileo of the error of the opinion and that he
abandon it; and that ``successively and without delay''
(\latin{successive ac incontinenti}), in the presence of the
cardinal and of witnesses, the Commissary of the Holy Office,
Michelangelo Seghizzi, O.P., enjoined on Galileo in the name of the
pope and the whole Congregation ``to relinquish altogether the
said opinion that the sun is the center of the world and immobile
and the earth moves, nor henceforth to hold, teach, or defend it
in any way whatever, verbally or in writing; otherwise the Holy
Office would proceed against him''; to which precept Galileo
``acquiesced and promised to obey.'' The minute names the
witnesses---members of Bellarmine's household---but bears no
signature of Galileo, of Bellarmine, or of a notary.

On Wednesday 3~March 1616 Bellarmine reported to the assembled
Congregation, the pope presiding, that ``Galileo the mathematician,
warned (\latin{monitus}) of the order of the Sacred Congregation
to abandon the opinion he had hitherto held\ldots\ acquiesced''
(\latin{acquievit}); and the pope ordered publication of the
Index decree (\work{Opere}~XIX, p.~278; N). Note what this
official report does not mention: no precept, no Seghizzi, no
threatened imprisonment---only the warning and the acquiescence.

The public act followed on 5~March 1616: the decree of the
Congregation of the Index, subscribed by Cardinal Sfondrati of
St.~Cecilia with the congregation's secretary---Bellarmine's name
appears nowhere on it---which declared ``that false Pythagorean doctrine, altogether contrary to
divine Scripture, of the mobility of the earth and the immobility
of the sun'' to be spreading; suspended Copernicus's \work{De
revolutionibus} and Zu\~niga's commentary on Job ``until they be
corrected'' (\latin{donec corrigantur}); and prohibited and
condemned Foscarini's book outright (\work{Opere}~XIX,
pp.~322--323, from the original printing; N; the English here
follows the public-domain Sturge translation of von Gebler,
checked against Favaro's text). The decree censures a doctrine as
``false and altogether opposed to Holy Scripture''; it does not
use the qualifiers' word ``heretical,'' it does not name Galileo,
and it does not touch his writings.

\subsection{26 May 1616: the certificate}

Rumor did what Rome had not: within weeks it was being said that
Galileo had abjured in the cardinal's hands and been given
penance. Galileo asked Bellarmine for a statement, and received
one written and signed in the cardinal's own hand. The autograph
survives in the Vatican dossier (car.\ 427r), where Galileo
produced it in his defense on 10~May 1633, together with a copy
in Galileo's own hand (car.\ 423r); Favaro prints both
(\work{Opere}~XIX, pp.~348 and 342; A at autograph level). It
reads, in project translation of Favaro's text:

\begin{studyframe}\small
``We, Roberto Cardinal Bellarmino, having heard that Signor
Galileo Galilei is calumniated or charged with having abjured in
our hand, and also with having been given salutary penances for
this, and being asked for the truth: we say that the said Signor
Galileo has not abjured in our hand, nor in the hand of any other
person here in Rome, nor, so far as we know, in any other place,
any opinion or doctrine of his, nor has he received salutary
penances or any other kind; but only there was announced to him
(\latin{gl'\`e stata denunziata}) the declaration made by Our Lord
[the Pope] and published by the Sacred Congregation of the Index,
in which it is contained that the doctrine attributed to
Copernicus---that the earth moves around the sun and that the sun
stands in the center of the world without moving from east to
west---is contrary to the Sacred Scriptures, and therefore may not
be defended or held. In witness whereof we have written and
subscribed the present with our own hand, this 26th day of May
1616. The same as above, Roberto Cardinal Bellarmino.''
\end{studyframe}

The certificate is Bellarmine's last documented act toward
Galileo; he died in 1621, eleven years before the trial. In 1633
it became the hinge of Galileo's defense---he testified that he
had relied on it, and its wording (``may not be defended or
held,'' with no ``taught'' and no \latin{quovis modo}) stood
visibly short of the unsigned minute's stricter formula, a
discrepancy the trial documents themselves discuss
(\work{Opere}~XIX; N).

\subsection{The 1616 dossier at a glance}

\begin{fourstable}{Date (1616)}{Document}{Witness status (per
Favaro)}{What it establishes}
\properrow{19 Feb}{Two propositions sent to the consultor
theologians}{dossier minute, car.\ 376v}{The examined propositions
in their exact wording}
\properrow{24 Feb}{Qualifiers' censure}{signed original, car.\
377r; eleven signatures, Bellarmine not among them}{The internal
theological censure (``formally heretical'' / ``at least erroneous
in faith''), never published as such}
\properrow{25 Feb}{Papal direction to Bellarmine}{Holy Office
minute, car.\ 378v}{The graduated order: warn; on refusal,
precept; on defiance, prison}
\properrow{26 Feb}{Warning and precept minute}{unsigned minute,
car.\ 378v--379r}{Warning by Bellarmine; immediate precept by
Seghizzi recorded; no signatures of Galileo, Bellarmine, or
notary}
\properrow{3 Mar}{Bellarmine's report to the Congregation}{session
record, dossier}{``Warned\ldots\ he acquiesced''; no mention of a
precept}
\properrow{5 Mar}{Index decree (public)}{original printing, car.\
380r; signed Sfondrati}{Copernicus and Zu\~niga suspended
\latin{donec corrigantur}; Foscarini condemned; ``heretical'' not
used; Galileo not named}
\properrow{26 May}{Bellarmine's certificate to Galileo}{autograph,
car.\ 427r; copy in Galileo's hand, car.\ 423r}{No abjuration, no
penance; the doctrine ``may not be defended or held''}
\end{fourstable}

\subsection{1633: the certificate in the trial}

Bellarmine was twelve years dead when the dossier reopened, but
his paper governed the endgame, and the 1633 records are the
certificate's reception history. In his first deposition (12
April 1633) Galileo produced a copy and quoted the Foscarini
letter from memory---``It appears to me that your Reverence and
Signor Galileo act wisely in contenting yourselves with speaking
\latin{ex suppositione} and not with certainty''---and testified
that Bellarmine had told him the opinion, ``taken absolutely,
must not be either held or defended,'' but might be held as a
supposition and so written about (von Gebler's account with the
deposition, trans.\ Sturge, pp.~203--204, from the dossier; N).
In his written defense of 10~May 1633 he submitted the autograph
itself, arguing that a certificate in which the operative words
were only ``defend'' and ``hold'' explained, without malice, his
failure to mention any stricter precept when seeking the
imprimatur (\work{Opere}~XIX, 348; the defense text at
car.~425--426; A/N). The sentence of 22~June 1633 turned the
shield around: it recited the certificate at length and then
ruled that it ``aggravated'' his case, since it stated the
doctrine contrary to Scripture, which he had nonetheless argued
as probable; and that the missing words \latin{docere} and
\latin{quovis modo} could not excuse silence about the precept
(\work{Opere}~XIX, 404--405; Sturge's translation, p.~233; N).
Whatever else the 1633 tribunal decided, on the documentary
point it conceded the fact this biography needs: the paper
Galileo held from Bellarmine really did say less than the
unsigned minute. The use each side made of that difference
belongs to Galileo's story, not Bellarmine's; it is noted here
because the certificate's afterlife is the largest single reason
his name remains in every history of science.

\subsection{The historiography, with ceilings}

The relation between the unsigned minute of 26~February and
everything else---the graduated papal order, Bellarmine's 3~March
report of a simple warning, and his certificate's ``only there
was announced to him''---has been the central problem of Galileo
scholarship since the Vatican dossier was opened to scholars in
the nineteenth century. The positions may be tabulated at
reported level: Emil Wohlwill and others argued the minute was a
later fabrication inserted to ground the 1633 prosecution; Karl
von Gebler, who collated the manuscript personally in 1877,
concluded against forgery on codicological grounds (the 1616
foliation sequence forbids the alleged insertion) while holding
it ``equally certain that the proceedings did not take place in
the rigid manner described in that annotation,'' the minute being
unsigned, un-notarized, and in tension with the graduated
instruction of 25~February (von Gebler, trans.\ Sturge, 1879,
pp.~337--338; K). Later scholarship has proposed many
mediations---that Seghizzi intervened prematurely before Galileo
had refused the warning; that the minute records honestly an
irregular act; that two accounts served two institutional
interests. This study asserts none of them.

\witnesslimit{Ceilings. (1) It is certain at document level that
Bellarmine warned Galileo and reported his acquiescence, and
certain that a minute records an immediate stricter precept whose
procedural regularity and exact course cannot now be verified.
(2) Bellarmine signed neither the qualifiers' censure nor the
Index decree; his documented personal acts in 1616 are the
warning, the report, and the certificate. (3) Nothing in 1616
tried, abjured, or punished Galileo, as the certificate itself
insists; the trial, abjuration, and condemnation belong to 1633,
after Bellarmine's death, and are outside this biography except
as reception. (4) None of this converts Bellarmine into either
the villain or the vindicated prophet of the affair: he shared
and enforced the theological consensus the 1616 acts express, and
he articulated, conditionally, the demonstration principle later
ages would quote against those acts. Both facts are in the
documents; their weighing is interpretation.}
